\documentclass[../main]{subfiles}

\begin{document}
\chapter*{Preface to the new typesetting}
\addcontentsline{toc}{chapter}{Preface To The New Typesetting}

The book was TeX'd up by the Texromancers, a latexing group. The credits for the typesetting of this book go to: Aareyan Manzoor, Abhishek Shivkumar, Yohan Wittgenstein, Kelvin Chan, RokettoJanpu, George Coote, John Cerkan, Mervyn Brumbach (III), Carl Sun, and others.  %put your names here.

Here is a link to a dyslexic friendly version: \url{https://aareyanmanzoor.github.io/assets/books/characteristic-classes-dyslexic.pdf} .

We added citations and references with hyperlinks. References to e.g. theorems/lemmas in the book are in blue, while citations to the bibliography is
in red. The bibliography also has URLs now, for easy access. Some of the books
in the bibliography had newer editions, so we went with those.

Some editor's notes were put in footnotes to point things out or talk about changes since the book came out.

A lot of little notation was changed to better fit the new latex'd version. For example, we changed tangent space to $\tangentspace{M}{x}$ rather than $D_x M$. Things like characteristic classes or homologies or special groups have formatting to emphasize them.

A section defining fibre bundles was added as the book uses the terminology in a few places. In the appendix, a section defining relative homology was added for the same reason.

Figures were redrawn to better fit the new style of the book.

The original published edition is Milnor and Stasheff's 1974 \emph{Characteristic Classes}. Corrections to this typesetting have been checked against the scanned Princeton edition available at \url{https://webhomes.maths.ed.ac.uk/~v1ranick/papers/milnstas.pdf}. They restore mathematical symbols, degrees, polynomial coefficients, omitted proof steps, diagram targets and arrows, and the intended internal and bibliographical references. The relative (co)homology notes now spell out the cochain extension required for the connecting homomorphism.

Corrections to errors already present in the original are distinguished by editor's footnotes. These include the Euler-characteristic hypotheses and references in Chapter~\ref{ch:2}, the universal-bundle homotopy target, a missing pullback star in the diagonal-class argument, connectedness in Problem~\ref{prob:13.C}, the map's homotopy class in Chapter~\ref{ch:18}, the final polynomial index and the coefficient index in Chapter~\ref{ch:19}, and the dimension and basis indices in Appendix~\ref{app:C}. Intentional notation modernization is retained.

To build the standard edition from the source directory, use a TeX installation with pdfLaTeX, a working Biber version compatible with the installed \texttt{biblatex}, and MakeIndex. Run \texttt{pdflatex main.tex}, \texttt{biber --output-safechars main}, \texttt{makeindex main.idx}, and then \texttt{pdflatex main.tex} twice (or use \texttt{latexmk -pdf main.tex}, which manages these dependencies). The repository's \texttt{.latexmkrc} configures Biber's \texttt{--output-safechars} option so that bibliographical accents, including combining marks, are emitted as pdfLaTeX-compatible commands. A single pdfLaTeX pass does not resolve the bibliography, index, or cross-references.

\end{document}